\chapter{2011年山东卷（文）}

\section{单选题}

\begin{problem}
设集合 \(M = \{x \mid (x + 3)(x - 2) < 0\}\)，\(N = \{x \mid 1 \leqslant x \leqslant 3\}\)，则 \(M \cap N =\)（\quad）

\choices
    {\([1,2)\)}
    {\([1,2]\)}
    {\((2,3]\)}
    {\([2,3]\)}
\end{problem}

\begin{answer}
A
\end{answer}

\begin{solution}
由 \((x+3)(x-2)<0\) 得 \(-3<x<2\)，所以 \(M=(-3,2)\). 又 \(N=[1,3]\)，故 \(M\cap N=[1,2)\)，选 A.
\end{solution}

\begin{problem}
复数 \(z = \frac{2 - \i}{2 + \i}\)（\(\i\) 为虚数单位）在复平面内对应的点所在的象限为（\quad）

\choices
    {第一象限}
    {第二象限}
    {第三象限}
    {第四象限}
\end{problem}

\begin{answer}
D
\end{answer}

\begin{solution}
分子、分母同乘 \(2-\i\)，得
\[
z=\frac{(2-\i)^2}{(2+\i)(2-\i)}=\frac{3-4\i}{5}=\frac35-\frac45\i
\]
实部为正、虚部为负，所以对应的点在第四象限，选 D.
\end{solution}

\begin{problem}
若点 \((a,9)\) 在函数 \(y = 3^{x}\) 的图象上，则 \(\tan \frac{a\pi}{6}\) 的值为（\quad）

\choices
    {0}
    {\(\frac{\sqrt{3}}{3}\)}
    {1}
    {\(\sqrt{3}\)}
\end{problem}

\begin{answer}
D
\end{answer}

\begin{solution}
由点 \((a,9)\) 在图象上，得 \(3^a=9=3^2\)，所以 \(a=2\). 因此
\[
\tan\frac{a\pi}{6}=\tan\frac{\pi}{3}=\sqrt{3}
\]
选 D.
\end{solution}

\begin{problem}
曲线 \( y = x^{3} + 11 \) 在点 \( P(1,12) \) 处的切线与 \( y \) 轴交点的纵坐标是（\quad）

\choices
    {\(-9\)}
    {\(-3\)}
    {9}
    {15}
\end{problem}

\begin{answer}
C
\end{answer}

\begin{solution}
设 \(f(x)=x^3+11\)，则 \(f'(x)=3x^2\)，所以曲线在 \(x=1\) 处的切线斜率为 \(f'(1)=3\). 切线方程为 \(y-12=3(x-1)\)，即 \(y=3x+9\). 令 \(x=0\)，得交点纵坐标为 \(9\)，选 C.
\end{solution}

\begin{problem}
已知 \(a\)，\(b\)，\(c \in \R\)，命题“若 \(a + b + c = 3\)，则 \(a^2 + b^2 + c^2 \geqslant 3\)”的否命题是（\quad）

\choices
    {若 \(a + b + c \neq 3\)，则 \(a^2 + b^2 + c^2 < 3\)}
    {若 \(a + b + c = 3\)，则 \(a^2 + b^2 + c^2 < 3\)}
    {若 \(a + b + c \neq 3\)，则 \(a^2 + b^2 + c^2 \geqslant 3\)}
    {若 \(a^2 + b^2 + c^2 \geqslant 3\)，则 \(a + b + c = 3\)}
\end{problem}

\begin{answer}
A
\end{answer}

\begin{solution}
设原命题为“若 \(p\)，则 \(q\)”，其中 \(p\) 为 \(a+b+c=3\)，\(q\) 为 \(a^2+b^2+c^2\geqslant3\). 否命题为“若非 \(p\)，则非 \(q\)”，即
\[
a+b+c\neq3\Longrightarrow a^2+b^2+c^2<3
\]
所以选 A.
\end{solution}

\begin{problem}
若函数 \(f(x) = \sin \omega x\)（\(\omega > 0\)） 在区间 \(\left[0, \frac{\pi}{3}\right]\) 上单调递增，在区间 \(\left[\frac{\pi}{3}, \frac{\pi}{2}\right]\) 上单调递减，则 \(\omega =\)（\quad）

\choices
    {\(\frac{2}{3}\)}
    {\(\frac{3}{2}\)}
    {2}
    {3}
\end{problem}

\begin{answer}
B
\end{answer}

\begin{solution}
函数在 \(x=\dfrac{\pi}{3}\) 的左侧递增、右侧递减，因此该点是局部最大值点. 由 \(f'(x)=\omega\cos(\omega x)\)，得
\(f'\left(\frac{\pi}{3}\right)=0\)，\(\cos\frac{\omega\pi}{3}=0\).
于是 \(\dfrac{\omega\pi}{3}=\dfrac{(2k+1)\pi}{2}\)，其中 \(k\) 为非负整数. 若 \(k\geqslant1\)，则第一个使 \(\cos(\omega x)=0\) 的正数点为 \(x=\dfrac{\pi}{2\omega}<\dfrac{\pi}{3}\)，其后 \(\cos(\omega x)<0\)，从而 \(f'(x)<0\)，与区间 \(\left[0,\dfrac{\pi}{3}\right]\) 上单调递增矛盾. 因此 \(k=0\)，所以 \(\omega=\dfrac32\)，选 B.
\end{solution}

\begin{problem}
设变量 \(x\)，\(y\) 满足约束条件 \(\begin{cases}
x + 2y - 5 \leqslant 0,\\
x - y - 2 \leqslant 0,\\
x \geqslant 0,
\end{cases}\) 则目标函数 \(z = 2x + 3y + 1\) 的最大值为（\quad）

\choices
    {11}
    {10}
    {9}
    {8.5}
\end{problem}

\begin{answer}
B
\end{answer}

\begin{solution}
约束条件等价于
\[
\begin{cases}
x\geqslant0,\\
y\leqslant\frac{5-x}{2},\\
y\geqslant x-2.
\end{cases}
\]
可行域的上边界为 \(y=\dfrac{5-x}{2}\)，两条边界相交于 \((3,1)\)，而 \(x=0\) 与上边界相交于 \(\left(0,\dfrac52\right)\). 目标函数 \(z=2x+3y+1\) 关于 \(y\) 递增，所以最大值在上边界上取得. 代入得
\[
z=2x+3\left(\frac{5-x}{2}\right)+1=\frac{x}{2}+\frac{17}{2}
\]
且可行域上边界的 \(x\) 满足 \(0\leqslant x\leqslant3\). 因此 \(z\) 在 \(x=3\)，\(y=1\) 时取得最大值 \(10\)，选 B.
\end{solution}

\begin{problem}
某产品的广告费用 \(x\) 与销售额 \(y\) 的统计数据如下表：

\begin{center}
\begin{tabular}{c|cccc}
广告费用 \(x\)（万元） & 4 & 2 & 3 & 5 \\
\hline
销售额 \(y\)（万元） & 49 & 26 & 39 & 54
\end{tabular}
\end{center}

根据上表可得回归方程 \(\hat{y} = \hat{b}x + \hat{a}\) 中的 \(\hat{b}\) 为 9.4，据此模型预报广告费用为 6 万元时销售额为（\quad）

\choices
    {\(63.6\text{ 万元}\)}
    {\(65.5\text{ 万元}\)}
    {\(67.7\text{ 万元}\)}
    {\(72.0\text{ 万元}\)}
\end{problem}

\begin{answer}
B
\end{answer}

\begin{solution}
由表中数据得 \(\bar{x}=\dfrac{4+2+3+5}{4}=\dfrac72\)，\(\bar{y}=\dfrac{49+26+39+54}{4}=42\).
回归直线经过样本中心 \((\bar{x},\bar{y})\)，所以
\[
\hat{a}=\bar{y}-\hat{b}\bar{x}=42-9.4\times3.5=9.1
\]
当广告费用 \(x=6\) 万元时，模型预测销售额为
\[
\hat{y}=9.4\times6+9.1=65.5\text{ 万元}
\]
故选 B.
\end{solution}

\begin{problem}
设 \(M(x_0, y_0)\) 为抛物线 \(C: x^2 = 8y\) 上一点，\(F\) 为抛物线 \(C\) 的焦点，以 \(F\) 为圆心，\(|FM|\) 为半径的圆和抛物线 \(C\) 的准线相交，则 \(y_0\) 的取值范围是（\quad）

\choices
    {\((0,2)\)}
    {\([0,2]\)}
    {\((2, + \infty)\)}
    {\([2, + \infty)\)}
\end{problem}

\begin{answer}
D
\end{answer}

\begin{solution}
抛物线 \(x^2=8y\) 的焦参数满足 \(4p=8\)，故 \(p=2\)，焦点为 \(F(0,2)\)，准线为 \(y=-2\). 由抛物线定义，点 \(M(x_0,y_0)\) 到焦点的距离等于到准线的距离，所以
\[
|FM|=y_0+2
\]
以 \(F\) 为圆心的圆与准线 \(y=-2\) 有公共点，当且仅当圆的半径不小于圆心到准线的距离 \(4\)，即
\[
y_0+2\geqslant4
\]
因此 \(y_0\geqslant2\)，取值范围为 \([2,+\infty)\)，选 D.
\end{solution}

\begin{problem}
函数 \( y = \frac{x}{2} - 2\sin x \) 的图象大致是（\quad）

\choices
    {\choicebitmap[width=0.15\paperwidth]{img/2011/shandong_liberal/q10_choice_a.png}}
    {\choicebitmap[width=0.15\paperwidth]{img/2011/shandong_liberal/q10_choice_b.png}}
    {\choicebitmap[width=0.15\paperwidth]{img/2011/shandong_liberal/q10_choice_c.png}}
    {\choicebitmap[width=0.15\paperwidth]{img/2011/shandong_liberal/q10_choice_d.png}}
\end{problem}

\begin{answer}
C
\end{answer}

\begin{solution}
函数 \(f(x)=\dfrac{x}{2}-2\sin x\) 满足 \(f(-x)=-f(x)\)，所以图象关于原点对称，且 \(f(0)=0\). 又
\(f'(x)=\frac12-2\cos x\)，\(f'(0)=-\frac32<0\).
令 \(\alpha=\arccos\dfrac14\)，则在 \([0,2\pi]\) 上，\(f'(x)=0\) 的两个点为 \(\alpha\) 和 \(2\pi-\alpha\)，函数依次递减、递增、递减. 因为 \(\sin\alpha=\dfrac{\sqrt{15}}{4}\)，且 \(\alpha<\dfrac{\pi}{2}<2<\sqrt{15}\)，所以
\(f(\alpha)=\frac{\alpha}{2}-2\sin\alpha<0\)，\(f(\pi)=\frac\pi2>0\)，\(f(2\pi)=\pi>0\).
因此图象从原点向右先下降到 \(x\) 轴下方，再上升到上方，随后下降但在 \(x=2\pi\) 处仍在 \(x\) 轴上方；左半部分由原点对称性确定. 只有选项 C 同时满足这些性质，故选 C.
\end{solution}

\begin{problem}
如图是长和宽分别相等的两个矩形. 给定下列三个命题：
\begin{circlelist}
\item 存在三棱柱，其正（主）视图，俯视图如图；
\item 存在四棱柱，其正（主）视图，俯视图如图；
\item 存在圆柱，其正（主）视图，俯视图如图.
\end{circlelist}
其中真命题的个数是（\quad）

\bitmapfigure[max width=6.0cm]{img/2011/shandong_liberal/q11_fig1.png}

\choices
    {3}
    {2}
    {1}
    {0}
\end{problem}

\begin{answer}
A
\end{answer}

\begin{solution}
设图中矩形的长、宽分别为 \(L\)、\(W\). 对于三棱柱，取棱柱的长度为 \(L\)，取其直三棱柱的三角形底面位于 \(yOz\) 平面内，并取该底面为两直角边均为 \(W\) 的直角三角形. 该三棱柱在 \(xz\) 平面上的正投影、在 \(xy\) 平面上的俯投影都分别覆盖 \(0\leqslant x\leqslant L\) 及一个长度为 \(W\) 的区间，故两个视图均为 \(L\times W\) 的矩形. 对于四棱柱，取长、宽、高分别为 \(L\)、\(W\)、\(W\) 的长方体，两个视图也均为 \(L\times W\) 的矩形. 对于圆柱，取轴长为 \(L\)、底面直径为 \(W\)，并使轴平行于矩形的长边，则正视图和俯视图均为 \(L\times W\) 的矩形. 因此三个命题都是真命题，真命题的个数为 \(3\)，选 A.
\end{solution}

\begin{problem}
设 \(A_{1}\)，\(A_{2}\)，\(A_{3}\)，\(A_{4}\) 是平面直角坐标系中两两不同的四点，若 \(\overrightarrow{A_{1}A_{3}} = \lambda \overrightarrow{A_{1}A_{2}}\) （\(\lambda \in \R\)），\(\overrightarrow{A_{1}A_{4}} = \mu \overrightarrow{A_{1}A_{2}}\) （\(\mu \in \R\)），且 \(\frac{1}{\lambda} + \frac{1}{\mu} = 2\)，则称 \(A_{3}\)，\(A_{4}\) 调和分割 \(A_{1}\)，\(A_{2}\). 已知点 \(C(c, 0)\)，\(D(d, 0)\) （\(c\)，\(d \in \R\)） 调和分割点 \(A(0, 0)\)，\(B(1, 0)\)，则下面说法正确的是（\quad）

\choices
    {\(C\) 可能是线段 \(AB\) 的中点}
    {\(D\) 可能是线段 \(AB\) 的中点}
    {\(C\)，\(D\) 可能同时在线段 \(AB\) 上}
    {\(C\)，\(D\) 不可能同时在线段 \(AB\) 的延长线上}
\end{problem}

\begin{answer}
D
\end{answer}

\begin{solution}
由定义，\(\overrightarrow{AC}=c\overrightarrow{AB}\)，\(\overrightarrow{AD}=d\overrightarrow{AB}\)，所以 \(c\neq0\)、\(d\neq0\)，且
\[
\frac1c+\frac1d=2\iff c+d=2cd
\]
若 \(C\) 是线段 \(AB\) 的中点，则 \(c=\dfrac12\)，代入上式得 \(\dfrac12+d=d\)，矛盾；同理 \(D\) 不可能是线段 \(AB\) 的中点. 另外，若 \(C\)，\(D\) 同时在线段 \(AB\) 上，则 \(0<c<1\)、\(0<d<1\). 当 \(0<c<\dfrac12\) 时，\(2c-1<0\)，由 \(d=\dfrac{c}{2c-1}\) 得 \(d<0\)；当 \(\dfrac12<c<1\) 时，\(0<2c-1<c\)，得 \(d>1\)，均不可能.

由 \(c+d=2cd\) 还可得
\[
(2c-1)(2d-1)=1
\]
若 \(C\)，\(D\) 同时在 \(AB\) 的延长线上，则两点或者分别位于 \(A\)、\(B\) 两侧，此时左式小于 \(0\)；或者同在一侧，此时两个因子同号且其绝对值都大于 \(1\)，左式大于 \(1\)，均不可能. 因而只有第四个说法正确，选 D.
\end{solution}

\section{填空题}

\begin{problem}
某高校甲，乙，丙，丁四个专业分别有 150，150，400，300 名学生，为了解学生的就业倾向，用分层抽样的方法从该校这四个专业共抽取 40 名学生进行调查，应在丙专业抽取的学生人数为\fillinblank{}.
\end{problem}

\begin{answer}
16
\end{answer}

\begin{solution}
四个专业共有 \(150+150+400+300=1000\) 名学生. 按比例分层抽样，丙专业应抽取
\[
40\times\frac{400}{1000}=16
\]
名学生.
\end{solution}

\begin{problem}
执行如图所示的程序框图，输入 \(l = 2\)，\(m = 3\)，\(n = 5\)，则输出的 \(y\) 的值是\fillinblank{}.

\texfigure[max width=6.8cm]{img_repaint/2011/shandong_liberal/q14_fig1.tex}
\end{problem}

\begin{answer}
68
\end{answer}

\begin{solution}
先计算 \(l^2+m^2+n^2=2^2+3^2+5^2=38\neq0\)，所以程序执行 \(y=70l+21m+15n\)，得到
\[
y=70\times2+21\times3+15\times5=278
\]
由于 \(278>105\)，第一次循环后 \(y=278-105=173\)；因 \(173>105\)，第二次循环后 \(y=173-105=68\). 此时 \(68\not>105\)，程序输出 \(y=68\).
\end{solution}

\begin{problem}
已知双曲线 \(\frac{x^2}{a^2} - \frac{y^2}{b^2} = 1\)（\(a > 0\)，\(b > 0\)）和椭圆 \(\frac{x^2}{16} + \frac{y^2}{9} = 1\) 有相同的焦点，且双曲线的离心率是椭圆离心率的两倍，则双曲线的方程为\fillinblank{}.
\end{problem}

\begin{answer}
\(\frac{x^2}{4} - \frac{y^2}{3} = 1\)
\end{answer}

\begin{solution}
椭圆的半长轴为 \(4\)，半短轴为 \(3\)，所以其焦距参数满足 \(c_E^2=4^2-3^2=7\)，离心率为 \(e_E=\dfrac{\sqrt7}{4}\). 双曲线与椭圆有相同焦点，故双曲线的焦距参数满足 \(c_H^2=7\). 又双曲线离心率为椭圆离心率的两倍，所以
\[
\frac{c_H}{a}=2e_E=\frac{\sqrt7}{2}
\]
从而 \(\dfrac{7}{a^2}=\dfrac74\)，得 \(a^2=4\). 由双曲线关系 \(c_H^2=a^2+b^2\)，得 \(b^2=7-4=3\). 因此双曲线方程为
\[
\frac{x^2}{4}-\frac{y^2}{3}=1
\]
\end{solution}

\begin{problem}
已知函数 \(f(x) = \log_a x + x - b\)（\(a > 0\)，且 \(a \neq 1\)）. 当 \(2 < a < 3 < b < 4\) 时，函数 \(f(x)\) 的零点 \(x_0 \in (n, n + 1)\)，\(n \in \N^*\)，则 \(n = \)\fillinblank{}.
\end{problem}

\begin{answer}
2
\end{answer}

\begin{solution}
由 \(a>2>1\)，函数在 \(x>0\) 上的导数为
\[
f'(x)=\frac{1}{x\ln a}+1>0
\]
所以 \(f(x)\) 在定义域上单调递增. 又因为 \(a>2\)，有 \(\log_a2<1\)，且 \(b>3\)，所以
\[
f(2)=\log_a2+2-b<1+2-b<0
\]
因为 \(a<3\)，有 \(\log_a3>1\)，且 \(b<4\)，所以
\[
f(3)=\log_a3+3-b>1+3-b>0
\]
由零点存在性及单调性，唯一零点满足 \(x_0\in(2,3)\)，故 \(n=2\).
\end{solution}

\section{解答题}

\begin{problem}
在 \(\triangle ABC\) 中，内角 \(A\)，\(B\)，\(C\) 的对边分别为 \(a\)，\(b\)，\(c\)，已知 \(\frac{\cos A - 2\cos C}{\cos B} = \frac{2c - a}{b}\).
\begin{enumerate}
\item 求 \( \frac{\sin C}{\sin A} \) 的值；
\item 若 \( \cos B = \frac{1}{4} \)，\( \triangle ABC \) 的周长为 5，求 \(b\) 的长.
\end{enumerate}
\end{problem}

\begin{solution}
\begin{enumerate}
\item 由题设分式有意义，\(\cos B\neq 0\). 根据正弦定理，\(a:b:c=\sin A:\sin B:\sin C\)，故
\[
\frac{2c-a}{b}=\frac{2\sin C-\sin A}{\sin B}
\]
将其代入题设等式，并利用 \(\sin B\neq 0\)、\(\cos B\neq 0\)，得
\[
\sin B(\cos A-2\cos C)=\cos B(2\sin C-\sin A)
\]
移项并按和角公式整理，得
\[
\sin B\cos A+\cos B\sin A=2(\sin B\cos C+\cos B\sin C)
\]
所以 \(\sin(A+B)=2\sin(B+C)\). 因为 \(A+B=\pi-C\)，\(B+C=\pi-A\)，所以 \(\sin C=2\sin A\). 再由正弦定理，
\[
\frac{\sin C}{\sin A}=2
\]

\item 由上面的结论及正弦定理，\(c=2a\). 由余弦定理和 \(\cos B=\frac14\)，得
\[
\frac14=\frac{a^{2}+c^{2}-b^{2}}{2ac}=\frac{5a^{2}-b^{2}}{4a^{2}}
\]
因此 \(b^{2}=4a^{2}\). 因边长为正，\(b=2a\). 又因三角形周长为 5，
\[
a+b+c=a+2a+2a=5a=5
\]
所以 \(a=1\)，从而 \(b=2\).
\end{enumerate}
\end{solution}

\begin{problem}
甲，乙两校各有 3 名教师报名支教. 其中甲校 2 男 1 女，乙校 1 男 2 女.
\begin{enumerate}
\item 若从甲校和乙校报名的教师中各任选 1 名，写出所有可能的结果，并求选出的 2 名教师性别相同的概率；
\item 若从报名的 6 名教师中任选 2 名，写出所有可能的结果，并求选出的 2 名教师来自同一学校的概率.
\end{enumerate}
\end{problem}

\begin{solution}
记甲校的两名男教师为 \(A_1\)，\(A_2\)，女教师为 \(A_3\)；记乙校的男教师为 \(B_1\)，两名女教师为 \(B_2\)，\(B_3\).

\begin{enumerate}
\item 从两校各选 1 名时，结果用有序对 \((A_i,B_j)\) 表示，所有可能结果为
\[
\begin{gathered}
(A_1,B_1),(A_1,B_2),(A_1,B_3),\\
(A_2,B_1),(A_2,B_2),(A_2,B_3),\\
(A_3,B_1),(A_3,B_2),(A_3,B_3).
\end{gathered}
\]
共有 \(3\times3=9\) 个等可能结果. 其中性别相同的结果有 \((A_1,B_1)\)，\((A_2,B_1)\)，\((A_3,B_2)\)，\((A_3,B_3)\)，共 4 个，因此选出的 2 名教师性别相同的概率为
\[
P_1=\frac{4}{9}
\]

\item 从 6 名教师中任选 2 名时，所有可能结果为
\[
\begin{gathered}
\{A_1,A_2\},\{A_1,A_3\},\{A_2,A_3\},\{B_1,B_2\},\{B_1,B_3\},\{B_2,B_3\},\\
\{A_1,B_1\},\{A_1,B_2\},\{A_1,B_3\},\{A_2,B_1\},\{A_2,B_2\},\{A_2,B_3\},\\
\{A_3,B_1\},\{A_3,B_2\},\{A_3,B_3\}.
\end{gathered}
\]
共有 \(\mathrm C_6^2=15\) 个等可能结果. 来自同一学校的结果有甲校内部的 \(\mathrm C_3^2=3\) 个和乙校内部的 \(\mathrm C_3^2=3\) 个，共 6 个. 因此选出的 2 名教师来自同一学校的概率为
\[
P_2=\frac{\mathrm C_3^2+\mathrm C_3^2}{\mathrm C_6^2}=\frac{6}{15}=\frac{2}{5}
\]
\end{enumerate}
\end{solution}

\begin{problem}
如图，在四棱台 \(ABCD - A_{1}B_{1}C_{1}D_{1}\) 中，\(D_{1}D \perp\) 平面 \(ABCD\)，底面 \(ABCD\) 是平行四边形，\(AB = 2AD\)，\(AD = A_{1}B_{1}\)，\(\angle BAD = 60^{\circ}\).
\begin{enumerate}
\item 证明： \(AA_{1} \perp BD\).
\item 证明：\(CC_{1} \parallel\) 平面 \(A_{1}BD\).
\end{enumerate}

\bitmapfigure[max width=6.0cm]{img/2011/shandong_liberal/q19_fig1.png}
\end{problem}

\begin{solution}
设 \(AD=s\)，取 \(A\) 为原点，\(AB\) 所在方向为 \(x\) 轴，\(AD\) 位于 \(xOy\) 平面内. 由 \(AB=2AD\) 和 \(\angle BAD=60^{\circ}\)，可设
\[
B=(2s,0,0)
\]
\[
D=\left(\frac{s}{2},\frac{\sqrt{3}s}{2},0\right)
\]
\[
C=\left(\frac{5s}{2},\frac{\sqrt{3}s}{2},0\right)
\]
设 \(DD_1=h\). 四棱台的两个底面平行，且 \(A_1B_1=AD=\frac12AB\)，所以两个底面对应长度的比为 \(\frac12\). 又因 \(D_1D\perp\) 平面 \(ABCD\)，有
\[
D_1=\left(\frac{s}{2},\frac{\sqrt{3}s}{2},h\right)
\]
\[
A_1=\left(\frac{s}{4},\frac{\sqrt{3}s}{4},h\right)
\]
\[
C_1=\left(\frac{3s}{2},\frac{\sqrt{3}s}{2},h\right)
\]

\begin{enumerate}
\item
\[
\overrightarrow{AA_1}=\left(\frac{s}{4},\frac{\sqrt{3}s}{4},h\right)
\]
\[
\overrightarrow{BD}=\left(-\frac{3s}{2},\frac{\sqrt{3}s}{2},0\right)
\]
因此
\[
\overrightarrow{AA_1}\cdot\overrightarrow{BD}=-\frac{3s^{2}}{8}+\frac{3s^{2}}{8}=0
\]
故 \(AA_1\perp BD\).

\item 有
\[
\overrightarrow{CC_1}=(-s,0,h)
\]
\[
\overrightarrow{DA_1}=\left(-\frac{s}{4},-\frac{\sqrt{3}s}{4},h\right)
\]
\[
\overrightarrow{DB}=\left(\frac{3s}{2},-\frac{\sqrt{3}s}{2},0\right)
\]
直接计算得
\[
\overrightarrow{DA_1}-\frac12\overrightarrow{DB}=(-s,0,h)=\overrightarrow{CC_1}
\]
向量 \(\overrightarrow{DA_1}\) 和 \(\overrightarrow{DB}\) 都在平面 \(A_1BD\) 内，所以它们的线性组合 \(\overrightarrow{CC_1}\) 是该平面内的方向向量.
事实上，若 \(C\) 在平面 \(A_1BD\) 内，则存在实数 \(\alpha\)，\(\beta\)，使 \(\overrightarrow{DC}=\alpha\overrightarrow{DB}+\beta\overrightarrow{DA_1}\). 比较第三个坐标得 \(\beta h=0\)，因 \(h>0\)，所以 \(\beta=0\)；再比较第二个坐标得 \(\alpha=0\)，这与 \(\overrightarrow{DC}=(2s,0,0)\) 矛盾.
故 \(C\notin\) 平面 \(A_1BD\)，从而 \(CC_1\parallel\) 平面 \(A_1BD\).
\end{enumerate}
\end{solution}

\begin{problem}
等比数列 \(\{a_{n}\}\) 中，\(a_{1}\)，\(a_{2}\)，\(a_{3}\) 分别是下表第一、二、三行中的某一个数，且 \(a_{1}\)，\(a_{2}\)，\(a_{3}\) 中的任何两个数不在下表的同一列.

\begin{center}
\begin{tabular}{c|ccc}
 & 第一列 & 第二列 & 第三列 \\
\hline
第一行 & 3 & 2 & 10 \\
第二行 & 6 & 4 & 14 \\
第三行 & 9 & 8 & 18
\end{tabular}
\end{center}

\begin{enumerate}
\item 求数列 \(\{a_{n}\}\) 的通项公式；
\item 若数列 \(\{b_{n}\}\) 满足：\(b_{n}=a_{n}+(-1)^{n}\ln a_{n}\)，求数列 \(\{b_{n}\}\) 的前 \(2n\) 项和 \(S_{2n}\).
\end{enumerate}
\end{problem}

\begin{solution}
由于 \(a_1\)，\(a_2\)，\(a_3\) 分别来自三行，且任意两项不在同一列，只需按三行选择三列的排列逐一检验.
\begin{enumerate}
\item 六种选择及等比条件 \(a_2^2=a_1a_3\) 的检验结果为
\[
\begin{gathered}
(3,4,18):\ 4^2\neq3\times18,\qquad (3,14,8):\ 14^2\neq3\times8,\\
(2,6,18):\ 6^2=2\times18,\qquad (2,14,9):\ 14^2\neq2\times9,\\
(10,6,8):\ 6^2\neq10\times8,\qquad (10,4,9):\ 4^2\neq10\times9.
\end{gathered}
\]
所以 \(a_1=2\)，\(a_2=6\)，\(a_3=18\)，公比为 \(q=3\)，数列通项为
\[
a_n=a_1q^{n-1}=2\cdot3^{n-1}
\]

\item 由 \(a_n=2\cdot3^{n-1}\)，有
\[
b_n=2\cdot3^{n-1}+(-1)^n\ln\left(2\cdot3^{n-1}\right)
\]
因此
\[
S_{2n}=\sum_{j=1}^{2n}2\cdot3^{j-1}+\sum_{j=1}^{2n}(-1)^j\bigl(\ln2+(j-1)\ln3\bigr)
\]
其中
\[
\sum_{j=1}^{2n}2\cdot3^{j-1}=\frac{2(3^{2n}-1)}{3-1}=3^{2n}-1
\]
并且
\[
\sum_{j=1}^{2n}(-1)^j=0
\]
以及
\[
\sum_{j=1}^{2n}(-1)^j(j-1)=\sum_{r=1}^{n}\bigl(-(2r-2)+(2r-1)\bigr)=n
\]
所以
\[
S_{2n}=3^{2n}-1+n\ln3
\]
\end{enumerate}
\end{solution}

\begin{problem}
某企业拟建造如图所示的容器（不计厚度，长度单位：米），其中容器的中间为圆柱形，左右两端均为半球形，按照设计要求容器的容积为 \(\frac{80\pi}{3}\) 立方米，且 \(l \geqslant 2r\). 假设该容器的建造费用仅与其表面积有关. 已知圆柱形部分每平方米建造费用为 3 千元，半球形部分每平方米建造费用为 \(c \) （\(c > 3\)）千元. 设该容器的建造费用为 \(y\) 千元.
\begin{enumerate}
\item 写出 \(y\) 关于 \(r\) 的函数表达式，并求该函数的定义域；
\item 求该容器的建造费用最小时的 \(r\).
\end{enumerate}

\bitmapfigure[max width=4.6cm]{img/2011/shandong_liberal/q21_fig1.png}
\end{problem}

\begin{solution}
\begin{enumerate}
\item 容器的容积等于圆柱的体积与两个半球的体积之和，因此
\[
\pi r^{2}l+\frac43\pi r^{3}=\frac{80\pi}{3}
\]
从而
\[
l=\frac{80}{3r^{2}}-\frac43r
\]
由 \(l\geqslant2r\) 得
\[
\frac{80}{3r^{2}}-\frac43r\geqslant2r\iff r^{3}\leqslant8
\]
又因半径 \(r>0\)，故 \(0<r\leqslant2\).

圆柱形部分的侧面积为 \(2\pi rl\)，两个半球的总表面积为 \(4\pi r^{2}\)，所以建造费用为
\[
y=3\cdot2\pi rl+c\cdot4\pi r^{2}=6\pi rl+4\pi cr^{2}
\]
代入 \(l\) 得
\[
y=\frac{160\pi}{r}+4\pi(c-2)r^{2}
\]
其中 \(0<r\leqslant2\).

\item 对 \(y\) 求导，得
\[
y'(r)=-\frac{160\pi}{r^{2}}+8\pi(c-2)r=\frac{8\pi}{r^{2}}\bigl((c-2)r^{3}-20\bigr)
\]
令 \(y'(r)=0\)，得到临界点
\[
r_0=\left(\frac{20}{c-2}\right)^{1/3}
\]
由于 \(c>3\)，有 \(c-2>0\)，且 \(y'(r)<0\)（当 \(0<r<r_0\)），\(y'(r)>0\)（当 \(r>r_0\)）.

当 \(3<c<\frac92\) 时，\(r_0>2\)，函数在定义域上递减，最小值在 \(r=2\) 时取得；当 \(c\geqslant\frac92\) 时，\(r_0\leqslant2\)，最小值在 \(r=r_0\) 时取得. 因此建造费用最小时
\[
r=\begin{cases}
2,&3<c<\dfrac92,\\
\left(\dfrac{20}{c-2}\right)^{1/3},&c\geqslant\dfrac92.
\end{cases}
\]
当 \(c=\frac92\) 时两种表达式均给出 \(r=2\).
\end{enumerate}
\end{solution}

\begin{problem}
在平面直角坐标系 \(xOy\) 中，已知椭圆 \(C: \frac{x^2}{3} + y^2 = 1\). 如图所示，斜率为 \(k \) （\(k > 0\)）且不过原点的直线 \(l\) 交椭圆 \(C\) 于 \(A\)，\(B\) 两点. 线段 \(AB\) 的中点为 \(E\). 射线 \(OE\) 交椭圆 \(C\) 于点 \(G\)，交直线 \(x = -3\) 于点 \(D(-3, m)\).
\begin{enumerate}
\item 求 \(m^{2} + k^{2}\) 的最小值；
\item 若 \(\left|OG\right|^{2}=\left|OD\right|\cdot\left|OE\right|\)，
\begin{enumerate}
\item 求证：直线 \(l\) 过定点；
\item 试问点 \(B\)，\(G\) 能否关于 \(x\) 轴对称？若能，求出此时 \(\triangle ABG\) 的外接圆方程. 若不能，请说明理由.
\end{enumerate}
\end{enumerate}

\bitmapfigure[max width=6.0cm]{img/2011/shandong_liberal/q22_fig1.png}
\end{problem}

\begin{solution}
设直线 \(l\) 的方程为 \(y=kx+b\)，其中 \(k>0\)，且 \(b\neq0\). 将其代入椭圆方程，得弦端点横坐标所满足的方程
\[
\left(k^{2}+\frac13\right)x^{2}+2kbx+b^{2}-1=0
\]
由韦达定理，线段 \(AB\) 的中点 \(E\) 的坐标为
\[
x_E=-\frac{3kb}{1+3k^{2}}
\]
\[
y_E=kx_E+b=\frac{b}{1+3k^{2}}
\]
即
\[
E=\frac{b}{1+3k^{2}}(-3k,1)
\]
\begin{enumerate}
\item 射线 \(OE\) 能与直线 \(x=-3\) 相交，且 \(k>0\)，所以 \(x_E<0\)，从而 \(b>0\). 直线 \(OE\) 的斜率为 \(-\frac{1}{3k}\)，故它与 \(x=-3\) 的交点为
\[
D\left(-3,\frac1k\right)
\]
所以 \(m=\frac1k\)，从而
\[
m^{2}+k^{2}=\frac1{k^{2}}+k^{2}\geqslant2
\]
当且仅当 \(k=1\) 时取等号，故最小值为 \(2\).

\item 在 \(\left|OG\right|^{2}=\left|OD\right|\cdot\left|OE\right|\) 的条件下，直线 \(OE\) 可参数表示为 \((x,y)=t(-3k,1)\). 将其代入椭圆方程，射线上的交点 \(G\) 对应
\[
t=\frac1{\sqrt{1+3k^{2}}}
\]
\[
G=\left(-\frac{3k}{\sqrt{1+3k^{2}}},\frac1{\sqrt{1+3k^{2}}}\right)
\]
于是
\[
\left|OG\right|^{2}=\frac{1+9k^{2}}{1+3k^{2}}
\]
又因为 \(b>0\)，
\[
\left|OE\right|=\frac{b\sqrt{1+9k^{2}}}{1+3k^{2}}
\]
\[
\left|OD\right|=\sqrt{9+\frac1{k^{2}}}=\frac{\sqrt{1+9k^{2}}}{k}
\]
将这些式子代入已知等式，得
\[
\frac{1+9k^{2}}{1+3k^{2}}=\frac{b(1+9k^{2})}{k(1+3k^{2})}
\]
所以 \(b=k\). 因此
\[
l:y=kx+k=k(x+1)
\]
故直线 \(l\) 恒过定点 \((-1,0)\).

\begin{enumerate}
\item 若点 \(B\)，\(G\) 关于 \(x\) 轴对称，则
\[
B=\left(-\frac{3k}{\sqrt{1+3k^{2}}},-\frac1{\sqrt{1+3k^{2}}}\right)
\]
将点 \(B\) 代入直线 \(y=k(x+1)\)，并记 \(s=\sqrt{1+3k^{2}}\)，得
\[
-\frac1s=k\left(1-\frac{3k}{s}\right)\iff 3k^{2}-ks=1
\]
于是 \(3k^{2}-1=ks>0\). 两边平方，得
\[
(3k^{2}-1)^{2}=k^{2}(1+3k^{2})\iff 6k^{4}-7k^{2}+1=0
\]
即 \((k^{2}-1)(6k^{2}-1)=0\). 由 \(3k^{2}-1>0\)，只能有 \(k=1\)，所以这种情形可以发生.

\item 当 \(k=1\) 时，\(l:y=x+1\)，且
\[
G=\left(-\frac32,\frac12\right)
\]
\[
B=\left(-\frac32,-\frac12\right)
\]
将直线 \(y=x+1\) 代入椭圆方程，得
\[
\frac{x^{2}}3+(x+1)^{2}=1\iff x(4x+6)=0
\]
所以另一个交点为 \(A=(0,1)\). 设三角形 \(ABG\) 的外接圆圆心为 \((u,v)\). 由于 \(BG\) 为竖直线段，其垂直平分线为 \(y=0\)，故 \(v=0\). 再由圆心到 \(A\)，\(B\) 的距离相等，
\[
u^{2}+1=\left(u+\frac32\right)^{2}+\frac14
\]
得 \(u=-\frac12\)，半径平方为 \(\frac54\). 因此外接圆方程为
\[
\left(x+\frac12\right)^{2}+y^{2}=\frac54
\]
即 \(x^{2}+y^{2}+x-1=0\).
\end{enumerate}
\end{enumerate}
\end{solution}
